\chapter{Implementação do Front-End}
\label{cap:frontend}

\section{Stack tecnológica}

O front-end utiliza React Native 0.81, Expo SDK 54, TypeScript 5.9, React Navigation 7, NativeWind 4 e React Native Reanimated \cite{react_native_docs,expo_docs,typescript_handbook,nativewind_docs,reanimated_docs}. Código em \texttt{cardly-rnative-app} \cite{cardly_frontend_repo}.

\section{Estrutura de pastas}

\begin{verbatim}
src/
├── api/           authApi, decksApi, dashboardApi, communityApi, adminApi
├── auth/          AuthContext.tsx
├── components/
│   ├── study/     FlipCard.tsx
│   └── ui/        ConfirmModal.tsx, toast.ts
├── navigation/    AuthStack, AppStack
├── screens/       Login, Register, Decks, Study, Dashboard, Admin...
└── theme/         colors.ts (paleta 2.0)
\end{verbatim}

Separação alinhada ao critério de avaliação sobre organização e responsabilidades \cite{clean_architecture_martin,react_hooks_docs}.

\section{Navegação e fluxos}

\subsection{Pilha de autenticação}

\texttt{AuthStack} contém \texttt{LoginScreen} e \texttt{RegisterScreen}. Após login bem-sucedido, \texttt{AuthContext} alterna para \texttt{AppStack}.

\subsection{Pilha principal}

\texttt{AppStack} agrupa:
\begin{itemize}
    \item \texttt{DashboardScreen}: hub com métricas e atalhos.
    \item \texttt{DecksScreen} / \texttt{DeckDetailScreen}: CRUD de assuntos e cartas.
    \item \texttt{StudySessionScreen}: sessão de estudo com flip e botões certo/errado/pular.
    \item \texttt{CommunityScreen}: assuntos públicos e clonagem.
    \item \texttt{AdminScreen}: visível apenas para \texttt{SUPERADMIN}.
\end{itemize}

React Navigation gerencia histórico e tipagem de parâmetros via \texttt{AppStackParamList} \cite{react_navigation}.

\section{Integração com API}

Cada módulo \texttt{*Api.ts} encapsula \texttt{fetch} \cite{mdn_fetch}:

\begin{enumerate}
    \item Monta URL base a partir de \texttt{EXPO\_PUBLIC\_API\_URL}.
    \item Injeta header \texttt{Authorization} quando autenticado.
    \item Serializa body JSON em POST/PUT/PATCH.
    \item Deserializa \texttt{PagedResponse} ou lança erro interpretado pela UI.
\end{enumerate}

Exemplo de fluxo no dashboard:

\begin{lstlisting}[language=Java,caption={Chamada tipada ao dashboard (TypeScript)},label={lst:dashboard-fetch}]
const data = await fetchDashboard(token);
setDashboard(data);
\end{lstlisting}

\section{Autenticação no cliente}

\texttt{AuthContext} provê \texttt{token}, \texttt{user}, \texttt{signIn}, \texttt{signOut} e \texttt{removeAccount}. Token persistido com \texttt{expo-secure-store} \cite{expo_secure_store}, atendendo recomendações OWASP para dados sensíveis em mobile \cite{owasp_mobile}.

Login Google envia ID token ao endpoint \texttt{/api/auth/google} \cite{google_oauth}.

\section{UI/UX e paleta de cores (padrão 2.0)}

Tokens definidos em \texttt{theme/colors.ts} e replicados no \texttt{tailwind.config.js} \cite{tailwind_docs}:

\begin{itemize}
    \item Primária: \texttt{\#2563EB} (ações principais).
    \item Fundo claro: \texttt{\#F8FAFC}; superfície: \texttt{\#FFFFFF}.
    \item Texto primário: \texttt{\#0F172A}; secundário: \texttt{\#64748B}.
    \item Feedback: success \texttt{\#22C55E}, warning \texttt{\#F59E0B}, error \texttt{\#EF4444}.
\end{itemize}

NativeWind aplica classes utilitárias (\texttt{className="bg-background-light"}) mantendo consistência visual \cite{nativewind_docs}.

\section{Componente FlipCard}

Requisito de UX: animação de virada ao toque (não hover), inspirada em flip cards CSS \cite{w3schools_flip}. Implementação com \texttt{Animated.spring} e rotação Y \cite{reanimated_docs}:

\begin{itemize}
    \item Estado \texttt{flipped} controlado pela tela de estudo.
    \item Frente exibe pergunta; verso exibe resposta.
    \item \texttt{backfaceVisibility: 'hidden'} evita artefatos visuais.
\end{itemize}

\section{Feedback ao usuário}

\subsection{Toasts}

\texttt{react-native-toast-message} exibe mensagens de sucesso e erro com moderação --- requisito de equilíbrio do professor \cite{toast_message_rn}. Helpers \texttt{showSuccess} e \texttt{showError} centralizam textos.

\subsection{Modais de confirmação}

Exclusões de conta, assuntos e cartas passam por \texttt{ConfirmModal}, reduzindo ações acidentais \cite{mobile_first_nielsen}.

\section{Tela de estudo}

\texttt{StudySessionScreen} orquestra:
\begin{enumerate}
    \item Carregamento de cartas vencidas (\texttt{dueOnly: true}).
    \item Exibição sequencial com \texttt{FlipCard}.
    \item Botões ``Errei'', ``Acertei'' e ``Pular'' mapeados a \texttt{AnswerCardRequest}.
    \item Toast e avanço automático para próxima carta.
\end{enumerate}

\section{Dashboard e calendário}

\texttt{DashboardScreen} exibe cards com \texttt{totalSubjects}, \texttt{totalCards}, \texttt{dueCards} e \texttt{answeredToday}. O calendário mensal com destaque de dias estudados está em implementação parcial (RF13); o design referencia aplicativos como Clue \cite{clue_app}.

\begin{figure}[htbp]
    \centering
    \fbox{\parbox{0.85\textwidth}{\centering
    \vspace{2cm}
    \textbf{[Figura 3 -- Print do Dashboard Cardly]}\\
    \vspace{2cm}
    }}
    \caption{Dashboard com indicadores de progresso (inserir captura de tela real)}
    \label{fig:dashboard-print}
\end{figure}

\section{Tela administrativa}

\texttt{AdminScreen} é renderizada condicionalmente quando \texttt{user.role === 'SUPERADMIN'}, atendendo requisito de front-end exclusivo para superadmin. Consome \texttt{adminApi} para buscas paginadas e formulários de criação/edição.

\section{Responsividade}

Layouts utilizam \texttt{flex-1}, padding consistente e tipografia escalável. Testes em emulador Android e dispositivo físico via Expo Go \cite{expo_docs,youtube_expo_typescript}.

\section{Hooks e estado}

Telas usam \texttt{useState}, \texttt{useEffect} e \texttt{useCallback} seguindo regras oficiais de Hooks \cite{react_hooks_docs}. Evita-se prop drilling excessivo delegando autenticação ao contexto global.

\section{Pendências de front-end}

\begin{itemize}
    \item Calendário visual completo no dashboard.
    \item Módulo de amizades (telas e integração API).
    \item Modo escuro completo com tokens \texttt{background.dark} (estrutura preparada na paleta).
\end{itemize}
